% chapters/anhang.tex
%
% A.1 dokumentiert die Reproduzierbarkeit der in Kapitel 4 beschriebenen
% Läufe (Labels app:reproduzierbarkeit, lst:prompt-glossar), A.2 ergänzende
% Ergebnistabellen (Tabellenkörper aus thesis/generated/).
% Benötigte Pakete in der Präambel: xcolor, tabularx, listings, booktabs.

\chapter{Anhang}
\label{chap:anhang}

\section{Reproduzierbarkeit der Läufe}
\label{app:reproduzierbarkeit}

Dieser Abschnitt ordnet den in Kapitel~\ref{chap:methodik} beschriebenen
Läufen Konfigurationsdateien, Codestände, Revisionen und
Ausgabeverzeichnisse zu und dokumentiert nachrangige Einstellungen.
Quellcode, Konfigurationsdateien, Auswertungsskripte und die aus den
Laufartefakten erzeugten Tabellen liegen im Repository
\url{https://github.com/luajosman/bachelor-thesis-llm-vs-encoder-ner}
(Commits nach Tabelle~\ref{tab:app-revisionen}; der Stand der
Zero-Shot-Läufe ist als Tag \texttt{zeroshot-runs-137bffd} markiert).
Konfigurationsdateien liegen dort unter \texttt{ba-ner/configs/}, die
Auswertung unter \texttt{thesis/scripts/aggregate\_results.py} mit den
Kennzahlen in \texttt{thesis/generated/summary.json}.
Die Laufartefakte selbst (Vorhersagen, Metriken, Adapter; rund 2~GB ohne
Modellgewichte) liegen relativ zum Ergebnisverzeichnis unter
\texttt{results/multinerd/} bzw. \texttt{results/crossner/} und werden
den Prüfenden auf Anfrage bereitgestellt.

\paragraph{Läufe und Artefakte.}
Tabelle~\ref{tab:app-laeufe} nennt für jede Konfiguration die
Konfigurationsdatei und das Ausgabeverzeichnis des Laufs mit Seed~42.
Für die Seeds 123 und 456 ersetzen die Konfigurationsdateien das Suffix
\texttt{\_canonical} durch \texttt{\_seed123} bzw. \texttt{\_seed456};
beim 27B-Modell wird es an \texttt{\_3ep} angehängt.
Die Ausgabeverzeichnisse erhalten das Suffix \texttt{-seed123} bzw.
\texttt{-seed456}.
Die Ausgabeverzeichnisse enthalten folgende Artefakte:
{\raggedright
\begin{itemize}
    \item alle Läufe: Testvorhersagen
    (\texttt{test\_\allowbreak predictions.json}) und Inferenzmetriken mit
    eingebetteten Laufmetadaten (\texttt{inference\_\allowbreak metrics.yaml});
    \item zusätzlich bei den trainierten Konfigurationen:
    Trainingsergebnisse (\texttt{results.yaml}), Trainingsverlauf
    (\texttt{training\_\allowbreak history.json}), ausgewählter Checkpoint
    (\texttt{best\_\allowbreak model} bzw.
    \texttt{best\_\allowbreak lora\_\allowbreak adapter}; Adapter 12,8, 42,5
    bzw. 637,6~MB), aufgelöste Konfiguration
    (\texttt{config\_\allowbreak resolved.yaml}), wissenschaftlicher
    Vertrag mit Code- und Revisions-Hashes
    (\texttt{scientific\_\allowbreak contract.yaml}), Paketversionen
    (\texttt{environment/}), Metadaten der Laufverwaltung
    (\texttt{run\_\allowbreak metadata.yaml}, \texttt{status.json};
    bei den 27B-Läufen zusätzlich \texttt{recovery-*.yaml}) sowie
    Evaluationsergebnisse einschließlich typbezogener Werte
    (\texttt{evaluation\_\allowbreak metrics.yaml}).
\end{itemize}\par}
Die CrossNER-Evaluationen verwenden die Checkpoints der Hauptstudie und
legen ihre Ergebnisse unter \texttt{results/crossner/} im Verzeichnis des
jeweiligen Laufs mit einem Unterordner je Domäne ab.

\begin{table}[ht]
  \centering
  \footnotesize
  \caption[Konfigurationsdateien, Ausgabeverzeichnisse und Codestände der Läufe]{Konfigurationsdateien, Ausgabeverzeichnisse (jeweils Seed~42)
  und Codestände der ausgewerteten Läufe; vollständige Commit-Hashes in
  Tabelle~\ref{tab:app-revisionen}.}
  \label{tab:app-laeufe}
  \setlength{\tabcolsep}{3pt}
  \begin{tabularx}{\linewidth}{@{}>{\raggedright\arraybackslash}p{3.3cm}%
      >{\raggedright\arraybackslash}p{4.35cm}>{\raggedright\arraybackslash}X%
      l@{}}
    \toprule
    \textbf{Konfiguration} & \textbf{Konfigurationsdatei} &
    \textbf{Ausgabeverzeichnis} & \textbf{Commit} \\
    \midrule
    DeBERTa-v3-base       & \texttt{deberta\_base\_canonical}  & \texttt{deberta-v3-base-canonical}  & \texttt{179eed8} \\
    DeBERTa-v3-large      & \texttt{deberta\_large\_canonical} & \texttt{deberta-v3-large-canonical} & \texttt{179eed8} \\
    Qwen3.5-0.8B QLoRA    & \texttt{qwen35\_08b\_canonical}    & \texttt{qwen35-08b-qlora-canonical} & \texttt{179eed8} \\
    Qwen3.5-4B QLoRA      & \texttt{qwen35\_4b\_canonical}     & \texttt{qwen35-4b-qlora-canonical}  & \texttt{179eed8} \\
    Qwen3.5-27B QLoRA     & \texttt{qwen35\_27b\_3ep}          & \texttt{qwen35-27b-qlora-3ep}       & \texttt{179eed8}$^{a}$ \\
    Qwen3.5-0.8B Zero-Shot & \texttt{qwen35\_08b\_zeroshot}$^{b}$ & \texttt{qwen35-08b-zeroshot}     & \texttt{137bffd} \\
    Qwen3.5-4B Zero-Shot  & \texttt{qwen35\_4b\_zeroshot}$^{b}$  & \texttt{qwen35-4b-zeroshot}      & \texttt{137bffd} \\
    Qwen3.5-27B Zero-Shot & \texttt{qwen35\_27b\_zeroshot}$^{b}$ & \texttt{qwen35-27b-zeroshot}     & \texttt{137bffd} \\
    CrossNER           & wie Hauptstudie                    & \texttt{results/crossner/}\dots      & \texttt{b620041} \\
    CrossNER, 1\,024 Tokens (Vorstudie) & wie Hauptstudie & \texttt{results/\allowbreak crossner\_output\_limit\_1024/\allowbreak smoke/}\dots & \texttt{ec49279} \\
    CrossNER, 1\,024 Tokens (alle Sätze, drei Seeds) & wie Hauptstudie & \texttt{results/\allowbreak crossner\_output\_limit\_1024/}\dots & \texttt{ff91342} \\
    Nachinferenz MultiNERD, 1\,024 Tokens (neun QLoRA-Läufe) & wie Hauptstudie & \texttt{results/\allowbreak multinerd\_output\_limit\_1024/}\dots & \texttt{ff91342} \\
    Greedy-Kontrolle Qwen3.5-0.8B & \texttt{qwen35\_08b\_zeroshot} & \texttt{qwen35-08b-zeroshot-greedy-v1} & \texttt{4679bf1}$^{c}$ \\
    Vorlauf DeBERTa-v3-base (Latenz auf A100) & \texttt{deberta\_base} & \texttt{deberta-v3-base} & \texttt{cdf10db}$^{c}$ \\
    \bottomrule
    \multicolumn{4}{@{}p{\linewidth}@{}}{\scriptsize $^{a}$ Start mit \texttt{179eed8};
      Fortsetzung nach den drei Unterbrechungen und Testinferenz mit
      \texttt{c01bcbe}.} \\
    \multicolumn{4}{@{}p{\linewidth}@{}}{\scriptsize $^{b}$ Dateistand von Commit
      \texttt{137bffd} (Tag \texttt{zeroshot-runs-137bffd}); der aktuelle
      Stand im Repository (Greedy-Decoding) weicht ab.} \\
    \multicolumn{4}{@{}p{\linewidth}@{}}{\scriptsize $^{c}$ Arbeitskopie mit lokalen,
      nicht erhaltenen Änderungen; Kontroll- bzw. Vorlauf außerhalb der
      Hauptauswertung.} \\
  \end{tabularx}
\end{table}

\paragraph{Ausführungszeiten und Cluster-Jobs.}
Die Trainings- und Inferenzläufe der Hauptstudie liefen im August 2026
auf dem DFKI-Cluster (Partitionen und Hauptspeicher in
Tabelle~\ref{tab:app-einstellungen}); ihre Slurm-Abrechnung liegt als
\texttt{results/slurm\_accounting\_2026-08.csv} vor, die 27B-Läufe
umfassen die Jobs 3254977, 3254980 und 3254983 (wiederholt eingereiht)
sowie 3267440, 3267443 und 3267446.
Die drei Zero-Shot-Läufe und der Greedy-Kontrolllauf entstanden am
1.~August bzw. 29.~Juli 2026 (27B auf einer H100~NVL, 0.8B und 4B auf
einer A100); ihre Job-IDs sind in den Laufartefakten nicht
protokolliert.
Die CrossNER-Evaluationen mit 256 Tokens liefen im Oktober 2026.
Die Nachinferenz mit 1\,024 Tokens lief am 28.~Oktober 2026 (Jobs
3673515 bis 3673523, H100~HBM3, Commit \texttt{ff91342}), der
vollständige CrossNER-Neulauf mit 1\,024 Tokens vom 28.\ bis
29.~Oktober 2026 (Jobs 3673475 bis 3673477, Aggregation 3673478); ein
erster Versuch (Jobs 3672273 bis 3672276) brach wegen eines
abweichenden Codestands vor der Inferenz ab.

\paragraph{Explorative Vorläufe.}
Tabelle~\ref{tab:app-vorlaeufe} dokumentiert die explorativen Vorläufe
(Abschnitt~\ref{sec:datenaufteilung}), je ein Lauf mit Seed~42, die auf
dem Testsplit ausgewertet wurden und nicht in die Hauptauswertung
eingehen; ihre Abweichung vom jeweiligen Seed-42-Lauf der Hauptstudie
beträgt höchstens 0,12 Punkte.

\begin{table}[ht]
  \centering
  \footnotesize
  \caption[Explorative Vorläufe]{Explorative Vorläufe (Seed~42) und ihr
  Mikro-F$_1$ auf dem MultiNERD-Testsplit (256 Ausgabetokens) gegenüber
  dem Seed-42-Lauf der Hauptstudie.}
  \label{tab:app-vorlaeufe}
  \setlength{\tabcolsep}{4pt}
  \begin{tabular}{llllrrr}
    \toprule
    \textbf{Vorlauf} & \textbf{Commit} & \textbf{GPU} & \textbf{Epochen} & \textbf{Mikro-F$_1$} &
      \textbf{Hauptstudie} & $\Delta$ \\
    \midrule
    DeBERTa-v3-base & \texttt{cdf10db}$^{a}$ & A100 & 5 & 95,58 & 95,49 & $+$0,10 \\
    DeBERTa-v3-large & \texttt{06b97be} & A100 & 4 & 95,84 & 95,83 & $+$0,01 \\
    Qwen3.5-0.8B QLoRA & \texttt{4679bf1} & A100 & 3 & 93,32 & 93,24 & $+$0,07 \\
    Qwen3.5-4B QLoRA & \texttt{137bffd} & A100 & 3 & 94,05 & 93,93 & $+$0,12 \\
    Qwen3.5-27B QLoRA & \texttt{137bffd}$^{a}$ & H100 NVL & 2 & 94,54 & 94,53 & $+$0,00 \\
    \bottomrule
    \multicolumn{7}{@{}l}{\scriptsize $^{a}$ Arbeitskopie mit lokalen, nicht erhaltenen \"Anderungen. Epochen: gew\"ahlte} \\
    \multicolumn{7}{@{}l}{\scriptsize Epoche (QLoRA) bzw. trainierte Epochen (Encoder); Hauptstudie: Seed-42-Lauf; $\Delta$ =} \\
    \multicolumn{7}{@{}l}{\scriptsize Vorlauf minus Hauptstudie.} \\
\end{tabular}

\end{table}

\begin{table}[ht]
  \centering
  \footnotesize
  \caption[Vollständige Commits und Revisionen]{Vollständige Commits und Revisionen.
  Die Modell- und Datensatzrevisionen gelten für die Trainingsläufe der
  Hauptstudie, den CrossNER-Versuch und die Sensitivitätsanalysen; die
  Zero-Shot-Läufe luden Modelle und Datensatz ohne explizite
  Revisionsangabe.}
  \label{tab:app-revisionen}
  \begin{tabular}{ll}
    \toprule
    \textbf{Komponente} & \textbf{Commit bzw. Revision} \\
    \midrule
    Code: Trainingsläufe der Hauptstudie     & \texttt{179eed8c68b49f51fff5498e411db22c88ef3c9a} \\
    Code: Fortsetzung und Inferenz 27B-QLoRA & \texttt{c01bcbe0f1be4356a2cf24b1bb501e03b7ab63be} \\
    Code: Zero-Shot-Läufe                    & \texttt{137bffdd9b4351b35a31e4550bde0c15b3c53a19} \\
    Code: CrossNER-Evaluationen              & \texttt{b620041ce75878603b37763831204f26fac7ab91} \\
    Code: CrossNER mit 1\,024 Tokens (Vorstudie) & \texttt{ec492799551cbfc9c47550e25006f11386f99ceb} \\
    Code: Sensitivitätsanalysen mit 1\,024 Tokens & \texttt{ff91342f1f206572c5e930e0be7d1929362ccb9c} \\
    Datensatz \texttt{Babelscape/multinerd}  & \texttt{2814b78e7af4b5a1f1886fe7ad49632de4d9dd25} \\
    CrossNER-Daten (Commit in \texttt{zliucr/CrossNER}) & \texttt{2e7ba2a7798c961e3f29fbc51252c5a8d40224bf} \\
    \texttt{Qwen/Qwen3.5-0.8B}               & \texttt{2fc06364715b967f1860aea9cf38778875588b17} \\
    \texttt{Qwen/Qwen3.5-4B}                 & \texttt{851bf6e806efd8d0a36b00ddf55e13ccb7b8cd0a} \\
    \texttt{Qwen/Qwen3.5-27B}                & \texttt{fc05daec18b0a78c049392ed2e771dde82bdf654} \\
    \texttt{microsoft/deberta-v3-base}       & \texttt{8ccc9b6f36199bec6961081d44eb72fb3f7353f3} \\
    \texttt{microsoft/deberta-v3-large}      & \texttt{64a8c8eab3e352a784c658aef62be1662607476f} \\
    \bottomrule
  \end{tabular}
\end{table}

\paragraph{Nachrangige Einstellungen.}
Tabelle~\ref{tab:app-einstellungen} ergänzt die Hyperparameter aus
Kapitel~\ref{chap:methodik} um Einstellungen, die das Ergebnis nur
mittelbar betreffen oder die Ausführung auf dem Cluster beschreiben.

\begin{table}[ht]
  \centering
  \footnotesize
  \caption[Nachrangige Trainings-, Auswertungs- und Ausführungseinstellungen]{Nachrangige Trainings-, Auswertungs- und
  Ausführungseinstellungen der trainierten Konfigurationen.}
  \label{tab:app-einstellungen}
  \setlength{\tabcolsep}{4pt}
  \begin{tabular}{lccccc}
    \toprule
    & \multicolumn{2}{c}{\textbf{DeBERTa-v3}} &
      \multicolumn{3}{c}{\textbf{Qwen3.5 (QLoRA)}} \\
    \textbf{Einstellung} & \textbf{base} & \textbf{large} &
      \textbf{0.8B} & \textbf{4B} & \textbf{27B} \\
    \midrule
    Batchgröße Validierung & 32 & 16 & 8 & 4 & 1 \\
    Generative Validierung & -- & -- &
      \multicolumn{3}{c}{200 Sätze, Batchgröße 1, greedy} \\
    Checkpoint-Intervall & Epoche & Epoche &
      \multicolumn{3}{c}{500 Schritte} \\
    Max.\ gespeicherte Checkpoints & \multicolumn{5}{c}{2} \\
    Logging-Intervall (Schritte) & 50 & 50 & 25 & 25 & 25 \\
    AdamW $\beta_1$; $\beta_2$; $\epsilon$ &
      \multicolumn{5}{c}{0,9; 0,999; $10^{-8}$} \\
    Gradient Clipping (max.\ Norm) & \multicolumn{5}{c}{1,0} \\
    Gradient Checkpointing & nein & nein &
      \multicolumn{3}{c}{ja (nicht-reentrant)} \\
    Dropout & \multicolumn{2}{c}{0,1 (Encoder und Kopf)} &
      \multicolumn{3}{c}{0,05 (LoRA)} \\
    Optimierungsschritte gesamt & 41\,025 & 32\,820 & 24\,615 & 24\,615 &
      24\,615 \\
    SLURM-Partition & H100 & H100 & H100 & A100-80GB & H100-Trails \\
    Angeforderter Hauptspeicher & 32\,GB & 48\,GB & 64\,GB & 96\,GB &
      160\,GB \\
    \bottomrule
  \end{tabular}
\end{table}

\paragraph{Zero-Shot-Prompt.}
Listing~\ref{lst:prompt-glossar} zeigt die Systemnachricht des
Glossar-Profils; Nutzernachricht und Chat-Template entsprechen
Listing~\ref{lst:prompt-task}.

\begin{lstlisting}[basicstyle=\ttfamily\scriptsize, breaklines=true,
  breakatwhitespace=true, columns=fullflexible, keepspaces=true,
  frame=single, captionpos=b, float=htbp,
  postbreak=\mbox{\textcolor{gray}{$\hookrightarrow$}\space},
  caption={Systemnachricht des im Zero-Shot-Regime verwendeten
  Glossar-Profils.},
  label={lst:prompt-glossar}]
You are an expert Named Entity Recognition (NER) system. Find every named-entity mention in the numbered token list. Named entities are specific named people, organizations, places, animals, biological entities, celestial bodies, diseases, events, foods, instruments, media or creative works, mythological entities, plants, times or dates, and vehicles. Do not label ordinary common nouns unless they are part of a specific named entity.
Use these exact MultiNERD labels: PER = Person; ORG = Organization; LOC = Location; ANIM = Animal; BIO = Biological entity; CEL = Celestial body; DIS = Disease; EVE = Event; FOOD = Food; INST = Instrument; MEDIA = Media or creative work; MYTH = Mythological entity; PLANT = Plant; TIME = Time or date; VEHI = Vehicle.
Return ONLY a JSON array. Do not use Markdown code blocks and do not add explanations. Extract all occurrences, including multi-token entities.
Each entity object must have exactly these fields: "start_token" (integer, inclusive), "end_token" (integer, exclusive), "text" (the exact text covered by the token span), and "type" (one of: PER, ORG, LOC, ANIM, BIO, CEL, DIS, EVE, FOOD, INST, MEDIA, MYTH, PLANT, TIME, VEHI).
Token indices must refer exactly to the numbered token list shown by the user. Check the entire token list before answering. Return [] only when none of the listed named-entity categories occurs.
\end{lstlisting}

\clearpage
\section{Ergänzende Ergebnistabellen}
\label{app:ergebnistabellen}

% Tabellenkörper werden von thesis/scripts/aggregate_results.py erzeugt
% (thesis/generated/); die Beispiele in tab:app-beispiele stammen aus den
% Testvorhersagen der genannten Läufe (Listenindex in test_predictions.json).

Tabelle~\ref{tab:app-seeds} nennt die Testwerte der einzelnen Seeds, aus
denen die Mittelwerte in Tabelle~\ref{tab:haupt} gebildet wurden, sowie
den Validierungs-F$_1$ und die Epoche des jeweils gewählten Checkpoints;
Tabelle~\ref{tab:app-limit-seeds} ergänzt die Nachinferenz mit
1\,024 Tokens (Tabelle~\ref{tab:limit-multinerd}) um die Werte je Seed;
Tabelle~\ref{tab:app-typen-prf} ergänzt Tabelle~\ref{tab:typen} um
Precision und Recall je Entitätstyp; die
Tabellen~\ref{tab:app-crossner-seeds} und~\ref{tab:app-crossner-typen}
ergänzen Tabelle~\ref{tab:crossner-erg} um die Werte der einzelnen Seeds
und um die Grobtypen je Domäne; Tabelle~\ref{tab:app-beispiele}
zeigt Beispiele typischer Fehler je Fehlerklasse aus
Abschnitt~\ref{sec:erg-fehlertypen}.

\begin{table}[ht]
  \centering
  \footnotesize
  \caption[Mikro-F$_1$ je Seed und Validierungs-F$_1$ der Checkpoints]{Mikro-F$_1$ (strikt, in Punkten) der einzelnen Seeds auf dem
  MultiNERD-Testsplit (Hauptläufe mit 256 Ausgabetokens) sowie
  Validierungs-F$_1$ des gewählten Checkpoints mit Epoche in Klammern
  (Encoder: vollständiger Validierungssplit; QLoRA: generative
  Validierung auf 200 Sätzen).}
  \label{tab:app-seeds}
  \begin{tabular}{llccc}
    \toprule
    \textbf{Modell} & \textbf{Regime} & \textbf{Seed 42} & \textbf{Seed 123} & \textbf{Seed 456} \\
    \midrule
    DeBERTa-v3-base & Fine-Tuning & 95,49 & 95,55 & 95,47 \\
    \quad Validierung (Epoche) & & 92,65 (4) & 92,57 (5) & 92,66 (5) \\
    DeBERTa-v3-large & Fine-Tuning & 95,83 & 95,78 & 95,84 \\
    \quad Validierung (Epoche) & & 93,08 (3) & 92,89 (4) & 93,05 (4) \\
    Qwen3.5-0.8B & QLoRA & 93,24 & 93,14 & 93,04 \\
    \quad Validierung (Epoche) & & 92,25 (3) & 93,15 (3) & 93,64 (3) \\
    Qwen3.5-4B & QLoRA & 93,93 & 94,02 & 93,96 \\
    \quad Validierung (Epoche) & & 95,45 (2) & 94,18 (3) & 95,05 (3) \\
    Qwen3.5-27B & QLoRA & 94,53 & 94,39 & 94,52 \\
    \quad Validierung (Epoche) & & 96,13 (3) & 94,68 (3) & 95,58 (3) \\
    \bottomrule
\end{tabular}

\end{table}

\begin{table}[ht]
  \centering
  \footnotesize
  \caption[Nachinferenz mit 1\,024 Tokens je Seed]{Nachinferenz der
  abgeschnittenen Antworten mit 1\,024 Ausgabetokens je Seed: Mikro-F$_1$
  bei 256 und 1\,024 Tokens, Differenz, Zahl der neu erzeugten und der
  weiterhin abgeschnittenen Antworten sowie Zahl der Antworten, deren
  erste 256 Tokens die gespeicherte Antwort reproduzieren.}
  \label{tab:app-limit-seeds}
  \setlength{\tabcolsep}{4pt}
  \begin{tabular}{llrrrrrr}
    \toprule
    \textbf{Modell} & \textbf{Seed} & \textbf{F$_1$ (256)} & \textbf{F$_1$ (1\,024)} & $\Delta$\textbf{F$_1$} &
      \textbf{\shortstack[c]{neu\\erzeugt}} & \textbf{\shortstack[c]{weiterhin\\abgeschnitten}} &
      \textbf{\shortstack[c]{Pr\"afix\\gleich}} \\
    \midrule
    Qwen3.5-0.8B & 42 & 93,24 & 94,35 & $+$1,11 & 68 & 0 & 67/68 \\
     & 123 & 93,14 & 94,29 & $+$1,14 & 68 & 0 & 67/68 \\
     & 456 & 93,04 & 94,16 & $+$1,12 & 67 & 0 & 67/67 \\
    \midrule
    Qwen3.5-4B & 42 & 93,93 & 95,06 & $+$1,14 & 72 & 0 & 70/72 \\
     & 123 & 94,02 & 95,15 & $+$1,14 & 67 & 0 & 66/67 \\
     & 456 & 93,96 & 95,11 & $+$1,15 & 67 & 0 & 66/67 \\
    \midrule
    Qwen3.5-27B & 42 & 94,53 & 95,68 & $+$1,15 & 65 & 0 & 63/65 \\
     & 123 & 94,39 & 95,57 & $+$1,18 & 68 & 0 & 67/68 \\
     & 456 & 94,52 & 95,67 & $+$1,15 & 62 & 0 & 61/62 \\
    \bottomrule
\end{tabular}

\end{table}

\begin{table}[ht]
  \centering
  \scriptsize
  \caption[Precision, Recall und F$_1$ je Entitätstyp]{Precision, Recall und F$_1$ je Entitätstyp (strikt, in Punkten,
  256 Ausgabetokens) für DeBERTa-v3-large und Qwen3.5-27B (QLoRA), Mittel
  über drei Seeds, sowie für Qwen3.5-27B im Zero-Shot-Regime (Einzellauf).
  Support nach Tabelle~\ref{tab:datensplits}.}
  \label{tab:app-typen-prf}
  \setlength{\tabcolsep}{3pt}
  \begin{tabular}{lrrrrrrrrrr}
    \toprule
    & & \multicolumn{3}{c}{\textbf{DeBERTa-v3-large}} &
      \multicolumn{3}{c}{\textbf{Qwen3.5-27B (QLoRA)}} &
      \multicolumn{3}{c}{\textbf{Qwen3.5-27B (Zero-Shot)}} \\
    \cmidrule(lr){3-5}\cmidrule(lr){6-8}\cmidrule(lr){9-11}
    \textbf{Typ} & \textbf{Support} & P & R & F$_1$ & P & R & F$_1$ & P & R & F$_1$ \\
    \midrule
    PER   & 5\,265 & 99,7 & 99,7 & 99,7 & 99,7 & 97,6 & 98,7 & 90,5 & 85,7 & 88,0 \\
    ORG   & 3\,309 & 98,6 & 99,0 & 98,8 & 99,0 & 98,8 & 98,9 & 87,7 & 76,8 & 81,9 \\
    LOC   & 12\,024 & 99,7 & 99,6 & 99,7 & 99,6 & 96,9 & 98,2 & 92,0 & 82,1 & 86,8 \\
    CEL   & 41 & 78,4 & 82,1 & 80,2 & 88,5 & 81,3 & 84,7 & 62,5 & 48,8 & 54,8 \\
    DIS   & 759 & 80,1 & 83,8 & 81,9 & 78,9 & 81,2 & 80,0 & 57,4 & 49,4 & 53,1 \\
    EVE   & 352 & 96,1 & 98,3 & 97,2 & 97,0 & 96,3 & 96,6 & 37,5 & 68,5 & 48,4 \\
    MEDIA & 458 & 97,4 & 98,3 & 97,8 & 97,4 & 98,3 & 97,8 & 62,9 & 41,5 & 50,0 \\
    ANIM  & 1\,604 & 76,7 & 82,3 & 79,4 & 74,4 & 73,6 & 74,0 & 54,0 & 38,7 & 45,0 \\
    BIO   & 8 & 72,6 & 87,5 & 79,0 & 59,2 & 79,2 & 67,5 & 1,6 & 100,0 & 3,2 \\
    FOOD  & 566 & 72,5 & 73,4 & 72,9 & 70,9 & 69,3 & 70,1 & 27,6 & 29,9 & 28,7 \\
    PLANT & 894 & 68,6 & 77,1 & 72,6 & 66,9 & 73,6 & 70,1 & 44,4 & 28,6 & 34,8 \\
    INST  & 12 & 85,5 & 80,6 & 82,7 & 97,2 & 97,2 & 97,2 & 4,9 & 33,3 & 8,6 \\
    MYTH  & 32 & 78,6 & 86,5 & 82,2 & 84,9 & 92,7 & 88,6 & 51,3 & 62,5 & 56,3 \\
    TIME  & 289 & 83,2 & 86,6 & 84,9 & 88,5 & 86,4 & 87,4 & 1,0 & 18,7 & 1,9 \\
    VEHI  & 32 & 90,0 & 92,7 & 91,3 & 92,0 & 93,8 & 92,8 & 10,0 & 65,6 & 17,4 \\
    \bottomrule
\end{tabular}

\end{table}

\begin{table}[ht]
  \centering
  \footnotesize
  \caption[CrossNER-Transfertest: F$_1$ je Domäne und Seed]{CrossNER-Transfertest
  mit 256 Ausgabetokens: F$_1$ (strikt, nur PER, ORG und LOC, in Punkten)
  je Domäne und Seed für DeBERTa-v3-large und Qwen3.5-27B (QLoRA).
  Gepoolt: Mikrowert über alle 2\,506 Sätze; Domänenmittel: ungewichtetes
  Mittel der fünf Domänen.}
  \label{tab:app-crossner-seeds}
  \setlength{\tabcolsep}{4pt}
  \begin{tabular}{lrrrrrrr}
    \toprule
    & & \multicolumn{3}{c}{\textbf{DeBERTa-v3-large}} &
      \multicolumn{3}{c}{\textbf{Qwen3.5-27B (QLoRA)}} \\
    \cmidrule(lr){3-5}\cmidrule(lr){6-8}
    \textbf{Dom\"ane} & \textbf{S\"atze} & \textbf{42} & \textbf{123} & \textbf{456} & \textbf{42} & \textbf{123} & \textbf{456} \\
    \midrule
    K\"unstliche Intelligenz & 431 & 82,14 & 81,48 & 82,32 & 75,16 & 74,40 & 74,87 \\
    Literatur & 416 & 83,28 & 84,07 & 84,65 & 77,21 & 78,67 & 76,96 \\
    Musik & 465 & 85,72 & 84,02 & 85,42 & 62,81 & 63,55 & 62,56 \\
    Politik & 651 & 86,79 & 86,18 & 87,81 & 71,33 & 72,44 & 70,99 \\
    Naturwissenschaften & 543 & 81,89 & 82,59 & 81,59 & 75,93 & 76,02 & 76,48 \\
    \midrule
    Gepoolt & 2\,506 & 84,98 & 84,52 & 85,46 & 71,49 & 72,27 & 71,31 \\
    Dom\"anenmittel & -- & 83,97 & 83,67 & 84,36 & 72,49 & 73,01 & 72,37 \\
    \bottomrule
\end{tabular}

\end{table}

\begin{table}[ht]
  \centering
  \footnotesize
  \caption[CrossNER-Transfertest: F$_1$ je Grobtyp und Domäne]{CrossNER-Transfertest
  mit 256 Ausgabetokens: F$_1$ je Grobtyp und Domäne (strikt, in Punkten,
  Mittel über drei Seeds) mit Support je Typ.}
  \label{tab:app-crossner-typen}
  \setlength{\tabcolsep}{3pt}
  \begin{tabular}{lrrrrrrrrr}
    \toprule
    & \multicolumn{3}{c}{\textbf{Support}} & \multicolumn{3}{c}{\textbf{DeBERTa-v3-large}} &
      \multicolumn{3}{c}{\textbf{Qwen3.5-27B (QLoRA)}} \\
    \cmidrule(lr){2-4}\cmidrule(lr){5-7}\cmidrule(lr){8-10}
    \textbf{Dom\"ane} & PER & ORG & LOC & PER & ORG & LOC & PER & ORG & LOC \\
    \midrule
    K\"unstliche Intelligenz & 227 & 173 & 83 & 86,6 & 72,5 & 92,8 & 76,5 & 68,0 & 85,8 \\
    Literatur & 742 & 110 & 200 & 87,1 & 69,9 & 82,1 & 80,6 & 67,3 & 73,0 \\
    Musik & 595 & 670 & 503 & 86,8 & 82,2 & 86,9 & 62,9 & 62,4 & 63,8 \\
    Politik & 839 & 1\,466 & 1\,017 & 91,4 & 86,7 & 84,0 & 72,4 & 75,7 & 64,5 \\
    Naturwissenschaften & 701 & 403 & 197 & 89,7 & 71,1 & 82,5 & 77,9 & 75,4 & 72,2 \\
    \bottomrule
\end{tabular}

\end{table}

\begin{table}[ht]
  \centering
  \footnotesize
  \caption[Beispiele typischer Fehler je Fehlerklasse]{Beispiele typischer Fehler je Fehlerklasse
  (Abschnitt~\ref{sec:erg-fehlertypen}) sowie der Formtreue-Kategorie
  \enquote{Unbekannter Typ} (Abschnitt~\ref{sec:erg-formtreue}) aus den
  Testvorhersagen mit 256 Ausgabetokens. Illustrative, nicht zufällig
  gezogene Auswahl; Nr.\ ist der nullbasierte Listenindex des Satzes in
  \texttt{test\_predictions.json}; Sätze aus dem MultiNERD-Testsplit
  \cite{tedeschi2022multinerd}. Referenz und Vorhersage als Span mit
  Typ; ein Strich bedeutet, dass die Referenz bzw. die Vorhersage keinen
  Span an dieser Stelle enthält. Seed~42, sofern nicht anders angegeben.}
  \label{tab:app-beispiele}
  \setlength{\tabcolsep}{3pt}
  \begin{tabularx}{\linewidth}{@{}>{\raggedleft\arraybackslash}p{0.85cm}
      >{\raggedright\arraybackslash}X
      >{\raggedright\arraybackslash}p{1.95cm}
      >{\raggedright\arraybackslash}p{1.95cm}
      >{\raggedright\arraybackslash}p{2.15cm}
      >{\raggedright\arraybackslash}p{2.55cm}@{}}
    \toprule
    \textbf{Nr.} & \textbf{Satz (gekürzt)} & \textbf{Referenz} & \textbf{Vorhersage} &
    \textbf{Fehlerklasse} & \textbf{Konfiguration} \\
    \midrule
    327 & He is also in the band Filthy Friends, playing bass and accordion. &
      Filthy Friends (PER) & Filthy Friends (ORG) & Typfehler &
      DeBERTa-v3-large und Qwen3.5-27B (QLoRA) \\
    396 & The lake is home to perch, pike-perch, bream, roaches, whitefishes,
      smelt and other species of fish. &
      perch (ANIM) & perch (FOOD) & Typfehler & DeBERTa-v3-large \\
    59 & In 2000, the entire campus was designated as an arboretum, the
      University of Georgia Campus Arboretum. &
      University of Georgia Campus Arboretum (ORG) &
      University of Georgia Campus (ORG); Arboretum (LOC) & Grenzfehler
      (mit Zusatz) & DeBERTa-v3-large \\
    59 & dito &
      University of Georgia Campus Arboretum (ORG) &
      University of Georgia (ORG); Campus Arboretum (LOC) &
      Grenz- und Typfehler & DeBERTa-v3-large (Seed~123) \\
    383 & The wilderness provides habitat for wild burro, desert tortoise,
      spotted bat and desert bighorn sheep. &
      burro (ANIM) & wild burro (ANIM) & Grenzfehler & Qwen3.5-27B (QLoRA) \\
    35 & It takes around 6,000 full-grown coconuts to produce one tonne of
      copra. &
      copra (FOOD) & -- & Auslassung & DeBERTa-v3-large \\
    4 & Within city limits, the southwestern fence lizard and whiptails are
      common. &
      whiptails (ANIM) & -- & Auslassung & Qwen3.5-27B (QLoRA) \\
    5 & Commonly seen mammals include the coyote, rock squirrel, Gunnison's
      prairie dog, desert cottontail, and black-tailed jackrabbit. &
      -- & desert cottontail (ANIM) & Zusatz &
      DeBERTa-v3-large und Qwen3.5-27B (QLoRA) \\
    10805 & Propranolol is a beta blocker that is mainly used to treat
      hypertension, but also prevents trembling. &
      -- & Propranolol (MED) & Unbekannter Typ & Qwen3.5-27B (Zero-Shot) \\
    2405 & The municipal holiday is Easter Monday. &
      Easter Monday (TIME) & Easter Monday (TME) & Unbekannter Typ &
      Qwen3.5-27B (Zero-Shot) \\
    \bottomrule
  \end{tabularx}
\end{table}
